\section{Evaluation Methodology}


To evaluate the impact of prompt engineering on algorithmic problem solving,
this study evaluates generated solutions using multiple quantitative metrics.

The evaluation focuses on solution correctness, algorithm quality, computational
efficiency, and execution performance.



\subsection{Correctness Evaluation}


The primary evaluation metric is solution correctness.


Each generated Python solution is executed against predefined test cases from
the benchmark dataset.


The correctness score is calculated as:


\[
Correctness = \frac{Passed\ Test\ Cases}{Total\ Test\ Cases}
\]


A higher correctness score indicates that the generated solution successfully
solves more programming tasks.



\subsection{Algorithm Selection Evaluation}


Beyond correctness, this research evaluates whether the generated solution
selects an appropriate algorithmic approach.


The evaluation compares:

\begin{itemize}

\item Selected algorithm pattern.

\item Expected optimal algorithm.

\item Data structure usage.

\end{itemize}


For example, a solution using a hash table for a two-sum problem is considered
more appropriate than a brute-force nested loop solution.



\subsection{Complexity Evaluation}


Generated solutions are analyzed based on computational complexity.


The evaluation considers:


\begin{itemize}

\item Time complexity.

\item Space complexity.

\end{itemize}


The generated complexity is compared with the expected optimal complexity
provided by the benchmark dataset.


A complexity score is assigned:


\[
Complexity\ Score =
\frac{Generated\ Efficiency}{Optimal\ Efficiency}
\]


This metric evaluates whether prompt engineering helps models produce more
optimized algorithms.



\subsection{Runtime Performance}


Each generated program is executed in a controlled environment.


Runtime performance is measured using:

\begin{itemize}

\item Execution time.

\item Memory usage.

\item Successful completion.

\end{itemize}


Runtime measurements are averaged across multiple test executions to reduce
random variation.



\subsection{Prompt Strategy Comparison}


The experiment compares three prompting strategies:


\begin{table}[h]

\centering

\begin{tabular}{lll}

\toprule

Strategy & Description & Purpose \\

\midrule

Baseline & Direct problem solving prompt & Measure default LLM capability \\

Optimized & Structured algorithm reasoning prompt & Evaluate prompt improvement \\

AI-Assisted & Multi-step engineering workflow & Evaluate collaborative workflow \\

\bottomrule

\end{tabular}


\caption{Prompt evaluation strategies}

\end{table}



The final comparison analyzes whether more structured prompting strategies
provide measurable improvements over direct LLM interaction.



\subsection{Statistical Analysis}


The experiment reports:

\begin{itemize}

\item Average correctness score.

\item Average runtime.

\item Complexity improvement.

\item Performance differences among prompt strategies.

\end{itemize}


These measurements are used to determine whether prompt engineering provides
significant benefits for algorithmic problem solving.
